\section{The reservation ledger is the real gate}
\label{sec:ledger}

Every term in \S\ref{sec:admission} is measured, and for a long time we believed
that settled the question. It does not, because \kn{START} is not the first gate.

The negotiator matches a job against the partitionable slot's advertised
\kn{Memory} attribute. That attribute is a \textbf{ledger}: configured capacity
minus what has already been promised to existing dynamic slots. It is not a
measurement, and a job is rejected before \kn{START} is ever evaluated if the
ledger says no.

\subsection{Measured cost}

On the sixteen-core, 31\,GB worker, running eight jobs that each \emph{requested}
3\,GB and used essentially nothing:

\begin{lstlisting}[style=term]
TotalMemory                    26942       %% configured, less RESERVED_MEMORY
8 dynamic slots x 3072      -  24576       %% promised
                              ------
advertised Memory               2366       %% machine refuses further work
MemAvailable (real)            24589       %% ten times as much genuinely free
\end{lstlisting}

Eight of sixteen cores idle. The policy was fully permissive throughout:
\kn{FITS\_MEM} evaluated \kn{3072 < 24589 - 7485}, which is true. The machine
declined work that its own policy would have admitted, because the matchmaker
never got far enough to ask.

\begin{trap}
This contradicts the rule the entire policy is built on. Admission was gated on
accounting after all, one layer above where we were looking.
\end{trap}

\subsection{Why HTCondor does this}

The ledger is not a defect in HTCondor. Its slot model is reservation-based:
when a dynamic slot is carved, its allocation is a guarantee held for that
slot's lifetime, whether the job touches it or not. Advertising measured free
memory would let two jobs match the same physical gigabytes and make the promise
to the first worthless.

That guarantee is real in a pool where \kn{request\_memory} is enforced. It is
\textbf{fictional here}, because \S\ref{sec:admission} disables enforcement
deliberately: nothing stops a job asking 3\,GB and using 30. The deployment was
therefore paying the full throughput cost of a reservation ledger while
receiving none of the protection it exists to provide.

\subsection{The attribute can be overridden}

HTCondor's documentation describes the two-address and slot-attribute model in
terms that suggest these values are the startd's to compute, and an earlier
revision of this note stated plainly that \kn{Cpus} and \kn{Memory} ``cannot
safely be overridden''. That was reasoned, never tested, and it is wrong.

A \kn{STARTD\_CRON} job publishing \kn{Memory} overrides the computed value.
The experiment carried a second attribute, \kn{MemOverrideProbe}, holding the
same number under a name the startd does not compute. Without that control, an
unchanged \kn{Memory} would have been ambiguous between \emph{the cron did not
run} and \emph{the attribute is protected}.

\begin{lstlisting}[style=term]
before      Memory 11582   TotalMemory 26942   MemOverrideProbe undefined
after       Memory 28122   TotalMemory 26942   MemOverrideProbe    28122
\end{lstlisting}

Both attributes carried the measured figure, and \kn{Memory} exceeded
\kn{TotalMemory}, so the startd had accepted a value it did not compute.

The second question is whether carving still works, and it does. Twelve jobs
requesting 3\,GB each, pinned to that machine, ran \textbf{concurrently}:

\begin{lstlisting}[style=term]
12 dynamic slots carved      36864 MB of nominal requests
TotalMemory                  26942 MB configured
claim refusals                   0
held jobs                        0
Cpus                         16 -> 4, decremented correctly
\end{lstlisting}

Under the ledger the same submission stops at eight.

\subsection{Publish headroom, not availability}

The value published is \kn{MemAvailable} minus a floor, not \kn{MemAvailable}
itself. The floor is
\kn{max(soft floor, RESERVED\_MEMORY)}.

Subtracting it is what keeps the advertisement agreeing with the policy.
Advertising the raw figure would let a job match and then be refused by
\kn{FITS\_MEM}, leaving it \st{Idle} against a machine that looks healthy, which
is precisely the failure shape of \S\ref{sec:silent}. Taking the maximum of the
two candidate floors is what lets one script be correct on every host: workers
reserve about ten percent, so the soft floor dominates, while a machine that
reserves more on purpose has its own reservation respected instead.

\subsection{Workers only}

The entry node keeps the ledger, and the script detects this itself rather than
being deployed selectively, publishing nothing when \kn{SCHEDD} appears in
\kn{DAEMON\_LIST}.

Its failure is not local: if the schedd cannot fork shadows, nothing anywhere in
the pool runs. The reservation also protects something real there, one shadow
per running job pool-wide, which unlike a job's \kn{request\_memory} is not a
guess. And its reservation is large by design, so a naive
\kn{MemAvailable} minus twenty-five percent would have advertised roughly double
what that machine is meant to offer. Detecting the role rather than configuring
it keeps one file correct fleet-wide.

\subsection{What this does not change}

No guard is weakened. Every memory term in \kn{START} and \kn{PREEMPT} reads the
cron's \kn{Host*} attributes, never the slot ad, so eviction behaves exactly as
it did when it was observed firing eleven seconds after pressure began. It
simply has more work to act on, and eviction becomes routine rather than rare.

Two side effects are worth stating. The published value lands on every slot, so
per-slot memory columns in \kn{condor\_status} stop being meaningful and only
the parent slot should be read. And \kn{SLOT\_WEIGHT} must stay at its \kn{Cpus}
default, or fair-share accounting would be computed from a number that no longer
describes an allocation.

Disk is deliberately left on the ledger. It is the one resource with no graceful
degradation.
